%% ==========================================================================
%%  MID-TERM TECHNICAL REPORT.  Companion to midterm.tex.
%%
%%  midterm.tex is the two-page per-member contribution summary the handout
%%  asks for, and is the graded artefact. THIS document is the technical
%%  record of what was actually built, and is supplementary.
%%
%%  Written incrementally, one section per working day, while the reasoning is
%%  still fresh. Reconstructing it on the last day is how the detail gets lost.
%%
%%  HOUSE STYLE (CLAUDE.md): no em dashes, no prose semicolons.
%%
%%  NOTE TO FUTURE EDITORS: write .tex files with the Write tool, not with a
%%  shell heredoc. Heredocs on this setup collapse every LaTeX line break from
%%  a double backslash to a single one, which silently destroys tables.
%%
%%  Build:  latexmk -pdf midterm_technical.tex
%% ==========================================================================
\documentclass[11pt]{article}
%% ==========================================================================
%%  preamble.tex --- packages, theorem environments, notation macros.
%%
%%  SHARED BY ALL THREE MILESTONE DOCUMENTS:
%%      survey.tex    Milestone 1, literature survey      (31 Aug 2026)
%%      midterm.tex   Milestone 2, per-member report      (12 Sep 2026)
%%      final.tex     Milestone 3, final report           ( 6 Nov 2026)
%%
%%  One preamble and one references.bib for the whole semester, so notation
%%  and citations cannot drift between documents.  If you need a symbol,
%%  add it HERE -- never redefine it inside a section file.
%% ==========================================================================

% ---------- page geometry & fonts -----------------------------------------
\usepackage[letterpaper,margin=1.15in]{geometry}
% Palatino text + matching math.  newpx is built on TeX Gyre Pagella, so it
% needs no URW base-35 fonts -- which a minimal TinyTeX install does not ship.
% If newpx is unavailable, `mathpazo' is the drop-in substitute wherever the
% URW fonts *are* present.
\usepackage[T1]{fontenc}
\usepackage[utf8]{inputenc}
\usepackage{newpxtext}
\usepackage{microtype}

% ---------- mathematics ---------------------------------------------------
% Order matters: amsmath must be loaded before newpxmath.
\usepackage{amsmath,amssymb,amsthm,mathtools}
\usepackage{newpxmath}
\usepackage{stmaryrd}          % \llbracket \rrbracket -- secret-sharing brackets

% ---------- floats, tables, algorithms ------------------------------------
\usepackage{booktabs}
\usepackage{multirow}
\usepackage{longtable}
\usepackage{array}
\usepackage{graphicx}
\graphicspath{{figures/}}
\usepackage{algorithm}
\usepackage[noend]{algpseudocode}
\floatname{algorithm}{Algorithm}
\usepackage{tikz}
\usetikzlibrary{arrows.meta,positioning,fit,backgrounds,calc}

% ---------- colour & links ------------------------------------------------
\usepackage{xcolor}
\definecolor{CiteGreen}{RGB}{0,120,60}
\definecolor{LinkRed}{RGB}{160,20,20}
\definecolor{UrlBlue}{RGB}{20,60,170}
\definecolor{DraftOrange}{RGB}{200,90,0}
\usepackage[
  colorlinks=true,
  citecolor=CiteGreen,
  linkcolor=LinkRed,
  urlcolor=UrlBlue,
  bookmarksnumbered=true
]{hyperref}
\usepackage{cleveref}

% ==========================================================================
%  Theorem environments.  One counter runs through
%  Definition/Theorem/Lemma/Proposition/Observation; Claim has its own.
% ==========================================================================
\theoremstyle{plain}
\newtheorem{theorem}{Theorem}[section]
\newtheorem{lemma}[theorem]{Lemma}
\newtheorem{proposition}[theorem]{Proposition}
\newtheorem{corollary}[theorem]{Corollary}
\newtheorem{definition}[theorem]{Definition}
\newtheorem{observation}[theorem]{Observation}
\newtheorem{conjecture}[theorem]{Conjecture}
\newtheorem{claim}{Claim}[section]

\theoremstyle{definition}
\newtheorem{example}[theorem]{Example}
\newtheorem{remark}[theorem]{Remark}
\theoremstyle{plain}

% NOTE.  cleveref cannot tell these apart --- they share the `theorem'
% counter, so \Cref would call every one of them "Theorem".  Write
% `Remark~\ref{...}', `Example~\ref{...}' explicitly.  \Cref is safe for
% sections, figures and tables.

% Bold, unpunctuated "Proof" head (amsthm defaults to italic "Proof.").
\makeatletter
\renewenvironment{proof}[1][\proofname]{\par
  \pushQED{\qed}%
  \normalfont \topsep6\p@\@plus6\p@\relax
  \trivlist
  \item[\hskip\labelsep\bfseries #1\@addpunct{}]\ignorespaces
}{%
  \popQED\endtrivlist\@endpefalse
}
\makeatother

% ==========================================================================
%  Draft notes.  Flip to \draftfalse for a clean submission copy; nothing
%  else in the document changes.
%
%      \todo{...}    something still to be written
%      \verify{...}  a claim that must be checked against the paper
%      \gap{...}     a place where the survey must argue the gap explicitly
%
%  Before submitting, find what is outstanding:
%      grep -rn '\\todo{\|\\verify{\|\\gap{' sections/ *.tex
% ==========================================================================
\newif\ifdraft
\drafttrue

\newcommand{\draftnote}[1]{%
  \ifdraft{\color{DraftOrange}\sffamily\small[\,#1\,]}\fi}
\newcommand{\todo}[1]{\draftnote{\textbf{TODO:} #1}}
\newcommand{\verify}[1]{\draftnote{\textbf{VERIFY:} #1}}
\newcommand{\gap}[1]{\draftnote{\textbf{GAP-ARGUMENT:} #1}}

% ==========================================================================
%  PROJECT NOTATION.
%
%  Fixed here, once, for all three documents.  Follows [NUDGE] where the two
%  base papers disagree, since its notation is the more standard.  Do not
%  introduce a competing symbol in a section file --- add it here instead,
%  and say so at the weekly sync so nobody duplicates it.
% ==========================================================================

% ---- parties and sharing -------------------------------------------------
\newcommand{\party}[1]{P_{#1}}                 % server P_0, P_1, P_2
\newcommand{\shr}[1]{\llbracket #1 \rrbracket} % replicated / additive sharing
\newcommand{\ring}{\mathcal{R}}                % the ring Z_{2^b}
\newcommand{\Ztwo}[1]{\mathbb{Z}_{2^{#1}}}
\newcommand{\bits}{b}                          % ring bit-length
\newcommand{\fracbits}{t}                      % fixed-point fractional bits
\newcommand{\seclen}{\lambda}                  % PRF seed / security parameter

% ---- the recommendation problem -----------------------------------------
\newcommand{\users}{m}                         % number of users
\newcommand{\items}{n}                         % number of items
\newcommand{\dime}{d}                          % embedding dimension
\newcommand{\iters}{\ell}                      % power-iteration rounds
\newcommand{\topk}{k}                          % recommendations returned
\newcommand{\ratings}{\mathbf{U}}              % user-item rating matrix, m x n
\newcommand{\userembed}{\mathbf{A}}            % user embeddings, m x d (SECRET)
\newcommand{\itemembed}{\mathbf{B}}            % item embeddings, d x n (PUBLIC)
\newcommand{\uvec}[1]{\mathbf{u}^{(#1)}}       % user i's rating vector
\newcommand{\avec}[1]{\mathbf{a}^{(#1)}}       % user i's embedding

% ---- protocols -----------------------------------------------------------
\DeclareMathOperator{\Trunc}{Trunc}
\DeclareMathOperator{\ApproxNormalize}{ApproxNormalize}
\DeclareMathOperator{\ApproxFactor}{ApproxFactor}
\DeclareMathOperator{\SetOrthogonal}{SetOrthogonal}
\DeclareMathOperator{\Normalize}{Normalize}
\DeclareMathOperator{\Gen}{Gen}
\DeclareMathOperator{\Eval}{Eval}
\DeclareMathOperator{\EvalFull}{EvalFull}
\newcommand{\dpf}{\mathrm{DPF}}
\newcommand{\fss}{\mathrm{FSS}}

% ---- shorthand for the two base papers ----------------------------------
%  Use \pirsona and \nudge in prose so the naming stays uniform; the macros
%  cite as well as name, so `\nudge{} publishes $\itemembed$ in the clear'
%  produces both the name and the citation.
\newcommand{\pirsona}{\textsc{Pirsona}\,\cite{pirsona2021}}
\newcommand{\nudge}{\textsc{Nudge}\,\cite{nudge2026}}
\newcommand{\pirsonaplain}{\textsc{Pirsona}}
\newcommand{\nudgeplain}{\textsc{Nudge}}
\newcommand{\sys}{\textsc{OblivRec}}           % our system, working name

% ---- generic -------------------------------------------------------------
\newcommand{\N}{\mathbb{N}}
\newcommand{\R}{\mathbb{R}}
\newcommand{\Z}{\mathbb{Z}}
\newcommand{\set}[1]{\left\{#1\right\}}
\newcommand{\abs}[1]{\left\lvert#1\right\rvert}
\newcommand{\norm}[1]{\left\lVert#1\right\rVert}
\newcommand{\inner}[2]{\left\langle #1, #2 \right\rangle}
\DeclareMathOperator*{\argmax}{arg\,max}
\DeclareMathOperator*{\argmin}{arg\,min}
\DeclareMathOperator{\poly}{poly}


\begin{document}

\title{\bfseries OblivRec: Private Serving and Delivery\\[0.3em]
       \large Mid-Term Technical Report, CS670}

\author{Mainak Sarkar (230619) \and Shrasti Dwivedi (230982) \and
        Aditya Anand (230065) \and Shravan Agrawal (230984)}

\date{CS670: Cryptographic Techniques for Privacy Preservation\\
      Indian Institute of Technology Kanpur \\
      12 September 2026}

\maketitle

\begin{abstract}
\noindent
This report records the first implementation slice of OblivRec, a
recommender that composes the three-party private training core of
\nudgeplain{} with the private delivery layer of \pirsonaplain{}. The slice
covers private serving and delivery against a model trained in the clear,
which is stage S1 of the staging fixed in the requirements. We report a
distributed point function written from the Boyle--Gilboa--Ishai
construction, a DPF-based private information retrieval read layer over the
MovieLens catalogue, a cleartext power-iteration factorisation used as the
quality oracle, and an experiment quantifying what an observer learns from
watching which items a user fetches. Private training itself is out of scope
here and is scheduled for the following phase.
\end{abstract}

\tableofcontents
\newpage

\section{Scope of this slice}
\label{sec:tech:scope}

The full system has two halves. Private \emph{training} learns a
factorisation of a rating matrix that no server may see, and private
\emph{delivery} fetches the recommended item without revealing which item it
was. The requirements document fixes the build order as S1, then S2, then
S3: serving and delivery first, then training, then the composition. S1 is
first because it is lower risk, it demonstrates something end to end early,
and the distributed point function it needs is a prerequisite for the
non-linear gates that private training will need later.

This report covers S1 only. The model is trained \emph{in the clear} by the
Python oracle of Section~\ref{sec:tech:oracle}, and everything downstream of
that model is private. Private training is the subject of the next phase and
no claim is made about it here.

\subsection{What is deliberately absent}
\label{sec:tech:scope:absent}

Three components specified in the design draft are not built in this slice,
and each was cut for a stated reason rather than for lack of time.

A distributed comparison function is not implemented. Its only consumers are
the truncation and normalisation protocols, both of which belong to private
training. The design draft asserts that a distributed point function and a
distributed comparison function are the same primitive. They are not, and the
reference implementation of \nudgeplain{} keeps them in separate modules,
which corroborates the distinction.

Oblivious top-$\topk$ selection is not implemented. The user reconstructs the
score vector and selects the top $\topk$ locally, so no server learns the
selection regardless of whether an oblivious selector exists. The apparatus
would protect something that is already protected.

Network emulation is not used. The development machine has neither Docker nor
\texttt{tc netem}, so every measurement in this report is taken on the
\texttt{local} profile and is labelled as such. Wide-area numbers are
deferred rather than estimated.

Batching is not implemented. \textsc{Architecture} \S8 gives the network layer
``framing, batching, \texttt{flush()}'', and the first two words are built. But
this slice exchanges a handful of frames per query, so there is nothing to
batch; \texttt{flush()} exists and is documented as the no-op it honestly is
under blocking sockets with no user-space buffer. An unused optimisation would
be harder to justify than its absence.

Wide-area behaviour is not measured, for the same reason network emulation is
not used. Every number here is on the \texttt{local} profile and says so.

\medskip

\noindent
An earlier version of this section said networking itself was not implemented
and named it as the largest open gap in the slice. That was true when written
and is no longer: the three servers are now three separate operating-system
processes talking over TCP, and Section~\ref{sec:tech:net} reports the
transcript and distinguisher experiment that the phase's exit criterion asks
for.

\section{The cleartext oracle, and how many power iterations are needed}
\label{sec:tech:oracle}

Every private component in this project is checked against a cleartext twin.
The factorisation oracle is therefore the first thing built, before any
cryptography, and it doubles as baseline B1, the speed-of-light reference
against which recommendation quality is measured.

\subsection{Construction}
\label{sec:tech:oracle:construction}

The oracle implements $\ApproxFactor$ as \nudgeplain{} states it, rather than
calling a singular value decomposition. This is deliberate. The private
implementation in the next phase must reproduce this procedure step by step
under secret sharing, so the reference has to match the \emph{structure} of
the protocol and not merely its output. Concretely, for each of $\dime$
components a random vector is repeatedly multiplied by
$\ratings^{\top}\ratings$, made orthogonal to the components already found,
and normalised. Each converged component is then revealed in the clear and
becomes a row of $\itemembed$.

Revealing each row before computing the next is what keeps the
orthogonalisation cheap, because it is then a Gram--Schmidt step against
public vectors. It is also the reason the item embedding matrix is public to
every server, which is the premise of the leakage analysis in this project.

A singular value decomposition is still computed, but as a \emph{cross-check}
rather than as the implementation. Power iteration on
$\ratings^{\top}\ratings$ converges to the top-$\dime$ right singular
subspace, so the largest principal angle between the two subspaces should
approach zero. That check costs two lines and it caught a genuine question on
the first day.

\subsection{Quality is flat in the iteration count}
\label{sec:tech:oracle:ell}

At the default of $\iters = 10$ inner rounds the subspace angle is
$9.66 \times 10^{-1}$ radians, which is nowhere near zero. Sweeping $\iters$
resolves what that means.

\begin{table}[h]
\centering
\small
\caption{MovieLens-100K, split \texttt{u1}, $\dime = 16$. The subspace angle
falls by five orders of magnitude while recommendation quality does not move.
Raw data in \texttt{bench/results/oracle\_ell\_sweep.jsonl}.}
\label{tab:tech:ell}
\begin{tabular}{@{}rrrr@{}}
\toprule
$\iters$ & subspace angle (rad) & nDCG@20 & Recall@20 \\
\midrule
10  & $9.659 \times 10^{-1}$ & 0.4536 & 0.4693 \\
25  & $4.170 \times 10^{-1}$ & 0.4545 & 0.4677 \\
50  & $1.834 \times 10^{-1}$ & 0.4525 & 0.4650 \\
100 & $6.966 \times 10^{-2}$ & 0.4523 & 0.4662 \\
200 & $3.604 \times 10^{-3}$ & 0.4526 & 0.4664 \\
400 & $6.652 \times 10^{-6}$ & 0.4526 & 0.4664 \\
\bottomrule
\end{tabular}
\end{table}

\begin{figure}[h]
\centering
\includegraphics[width=0.68\linewidth]{figures/fig_ell}
\caption{The same data as Table~\ref{tab:tech:ell}, plotted. Note the axes:
nDCG@20 spans a range of 0.002 while the subspace angle spans five orders of
magnitude. Since $\iters$ multiplies the only interactive cost in private
training, the flat curve is what makes a fortyfold communication saving free.
Regenerated by \texttt{mingw32-make figures}.}
\label{fig:ell}
\end{figure}

Two conclusions follow, and they are separate.

The first is that the implementation is correct. The angle falls
monotonically to $6.7 \times 10^{-6}$ radians, so the procedure does converge
to the right subspace. The large angle at $\iters = 10$ is slow convergence
rather than a defect. The spectrum explains the rate: the ratio of the
sixteenth to the seventeenth squared singular value is $1.032$, a gap of three
percent, and power iteration converges as a power of that ratio. The first
component is easy, with a ratio of $6.44$. The sixteenth is not.

This distinction was nearly lost. A failing cross-check invites the reaction
of removing the cross-check, which would have converted a correct
implementation into a silently wrong one.

The second conclusion is the useful one. Across a fortyfold increase in
$\iters$, nDCG@20 stays between $0.4523$ and $0.4545$, a spread of about half
a percent with no monotone trend. Ranking does not require the singular
vectors, only a subspace good enough that the ordering of scores is stable.

\subsection{Why this matters for private training}
\label{sec:tech:oracle:consequence}

Under two-out-of-three replicated sharing the matrix--vector products are
non-interactive and therefore free. Truncation and normalisation are the
entire communication cost of training, and they run once per inner iteration,
so communication is linear in $\iters$.

The measurement therefore says that private training can run at $\iters = 10$
and the fortyfold saving in communication costs nothing in recommendation
quality. That is a design decision resting on evidence rather than on the
default value in the design draft, and it also discharges an instruction the
draft had already given and which had not been acted on, namely to measure
the residual against the cleartext oracle and report iterations against
quality.

The caveats are stated plainly. This is one dataset, one split, one seed and
one embedding dimension. The flatness may not survive at MovieLens-1M or at
$\dime = 32$, and it will be re-measured before the final report relies on it.
The metric uses binary relevance at a rating of four or above, so these
numbers are a within-project baseline and are not directly comparable to the
figures \nudgeplain{} reports on Netflix data under its own convention.

\section{The distributed point function}
\label{sec:tech:dpf}

A point function $f_{\alpha,\beta}$ takes the value $\beta$ at $x = \alpha$ and
zero everywhere else. A $(2,2)$-distributed point function splits it into two
keys, each of which individually hides both $\alpha$ and $\beta$, such that the
two evaluations recombine to $f_{\alpha,\beta}(x)$. The construction is that of
Boyle, Gilboa and Ishai, and it is written by us rather than imported, because
it is the component a viva will probe hardest and because the same
implementation serves both halves of this project: the private read layer in
Section~\ref{sec:tech:pir}, and the function secret sharing gates that private
training will need in the next phase.

\subsection{Why the keys are small}
\label{sec:tech:dpf:size}

The construction is a binary tree of pseudorandom seeds. Each key holds one
initial seed and one correction word per level, so key size is
\emph{logarithmic} in the domain while a naive one-hot query is linear. Table
\ref{tab:tech:keysize} gives measured sizes from the serialiser rather than
from the closed form, since the closed form is what the design draft got wrong.

\begin{table}[h]
\centering
\small
\caption{Measured DPF key size against the naive alternative. The packed
format is $1 + 4 + 16 + 18\,\mathrm{db} + \lceil \bits/8 \rceil$ bytes, and
\texttt{SizeBytes()} is asserted equal to the actual serialised length.}
\label{tab:tech:keysize}
\begin{tabular}{@{}rrrrrr@{}}
\toprule
domain bits & $\items$ & key, $\bits=64$ & key, $\bits=128$ & naive one-hot & compression \\
\midrule
10 & 1{,}024  & 209 B & 217 B & 8{,}192 B   & $39\times$ \\
11 & 2{,}048  & 227 B & 235 B & 16{,}384 B  & $72\times$ \\
12 & 4{,}096  & 245 B & 253 B & 32{,}768 B  & $134\times$ \\
16 & 65{,}536 & 317 B & 325 B & 524{,}288 B & $1654\times$ \\
\bottomrule
\end{tabular}
\end{table}

The design draft claims roughly 260 bytes at $\items = 4096$ against 32~KB
naive. The 32~KB figure is right. The 260-byte figure is not: the measured
value is \textbf{245 bytes}, and the discrepancy comes from the draft's formula
omitting both the initial seed and the self-describing header, while
simultaneously evaluating at 16 domain bits rather than the 12 that
$\items = 4096$ implies. The number in this report is measured by a test that
prints it, which is why the error surfaced.

\subsection{The sign convention}
\label{sec:tech:dpf:sign}

Textbook presentations of this construction define
$\Eval_b = (-1)^b\bigl(\mathrm{Convert}(s) + t \cdot \mathrm{CW}_{\text{last}}\bigr)$,
which yields the \emph{sum} convention $\Eval_0(x) + \Eval_1(x) = f(x)$. The
specification for this project instead states the invariant as a
\emph{difference}. We therefore omit the $(-1)^b$ factor and return
$\mathrm{Convert}(s) + t \cdot \mathrm{CW}_{\text{last}}$ directly, giving
%
\[
  \Eval(k_0, x) - \Eval(k_1, x) =
  \begin{cases}
    \beta & x = \alpha \\
    0     & \text{otherwise.}
  \end{cases}
\]
%
The two conventions are algebraically identical and differ only in where the
negation is placed. The reason to record the choice explicitly is that getting
it backwards makes every test fail in a way that resembles a tree traversal
bug rather than a sign error, which is an expensive place to spend a day.

\subsection{Testing: the structural invariant, checked at every node}
\label{sec:tech:dpf:testing}

A black-box test reports only that a leaf is wrong. The characteristic failure
of this construction is a mistaken control-bit update, which stays invisible
until enough levels have accumulated for it to change an output, and therefore
presents as ``correct at two domain bits, wrong at eight''. Bisecting that from
leaf values is slow.

The test suite therefore walks \emph{both} parties' trees in lockstep and
asserts the underlying structural invariant at every internal node. Writing
$(s_b, t_b)$ for party $b$'s state:

\begin{center}
\begin{tabular}{@{}ll@{}}
\toprule
On the path to $\alpha$  & $s_0 \neq s_1$ and $t_0 \oplus t_1 = 1$ \\
Off the path             & $s_0 = s_1$ and $t_0 \oplus t_1 = 0$ \\
\bottomrule
\end{tabular}
\end{center}

Off the path the two parties hold identical state, which is precisely why
their contributions cancel. On the path they differ and their control bits
disagree, which is what lets the final correction word inject $\beta$ at
exactly one leaf. Every correctness property of the scheme follows from these
two lines, so a violation names the exact level and node at which the
construction first went wrong.

This checker was written \emph{before} the generator was trusted, and both
were exercised together on the first run.

\subsection{Correctness results}
\label{sec:tech:dpf:results}

The structural invariant holds at every node for every $\alpha$ in every
domain up to eight bits, and at sampled $\alpha$ up to eleven bits.

The difference invariant is checked against a brute-force point function
oracle for every $\alpha$ and every $x$: exhaustively to ten domain bits on
the 64-bit ring and nine on the 128-bit ring, then for all $x$ at sampled
$\alpha$ at eleven and twelve bits. Edge cases that off-by-one errors land on
are covered explicitly, namely $\alpha$ at the first and last index, $\beta = 0$
where everything must cancel, and $\beta = -1$ where the ring wraps. A separate
check confirms that $\beta = 1$ reconstructs exactly, since the read layer
depends on that with no tolerance for an error term.

Traversal is implemented once and shared by single-point evaluation,
whole-domain evaluation and the test checker, so there is one traversal
implementation to be wrong rather than three.

\subsection{What is implemented but not yet verified}
\label{sec:tech:dpf:pending}

Whole-domain evaluation and key serialisation are implemented, because the
header is a frozen cross-component contract and had to be complete, but they
are so far exercised only indirectly. Extending the exhaustive invariant to
sixteen domain bits through whole-domain evaluation, asserting that it agrees
with single-point evaluation at every index, and round-tripping serialised keys
are the next steps. The key sizes in
Table~\ref{tab:tech:keysize} do come from the real serialiser, so that path is
at least exercised.

\section{Private serving and private delivery}
\label{sec:tech:pir}

This section covers the half of the system that is built and running: the user
is scored against the catalogue without any server learning the user's
embedding, and the resulting records are fetched without either serving party
learning which records were fetched. Private \emph{training} is the next phase
and is not claimed here.

\subsection{The catalogue is fixed-width, and that is a security property}
\label{sec:tech:pir:catalogue}

A private read must return the same number of bytes whatever it reads, or the
response length itself leaks the record. MovieLens-100K's \texttt{u.item} is
variable-length text, so it is packed into a fixed record of $L = 256$ bytes.

Two fields are dropped. Field~3, the video release date, is \emph{empty on all
1682 lines} --- measured, not assumed --- and field~4, the IMDb URL, is the
largest field and is used by nothing downstream. What remains is
\texttt{id\,|\,title\,|\,date\,|\,19 genre flags}, whose longest instance is
117 bytes of content, leaving 139 bytes of headroom inside $L$. Every record is
checked against that budget at load time, so a future dataset with a longer
title fails loudly at startup rather than truncating silently at answer time.

Fields are carried with explicit one-byte length prefixes rather than a
delimiter. Nine titles contain Latin-1 bytes and the file does not decode as
UTF-8 at all, so it is read as bytes and never re-encoded; with arbitrary bytes
in the payload, any delimiter choice risks colliding with content, whereas a
length prefix cannot. Three records break a naive parser and each is handled
explicitly: id~267 is \texttt{unknown} with an empty date and empty URL,
id~1242 is the width driver, and one release date is \texttt{4-Feb-1971} rather
than the \texttt{DD-Mon-YYYY} that every other date follows --- which never
mattered, because dates are carried as opaque text and never parsed into a
structured type.

The 1682 items pad to a domain of $2048$, so $\mathrm{db} = 11$. Padding
records are all-zero. This leaks nothing precisely because record width is
fixed regardless of content: a fetch of a padding slot is indistinguishable
from a fetch of a real one.

\subsection{The private read}
\label{sec:tech:pir:read}

The client wants record $\alpha$ without either server learning $\alpha$. It
generates a DPF key pair for the point function that is $1$ at $\alpha$, sends
one key to each of $\party{0}$ and $\party{1}$, and each server expands its key
over the whole domain and takes an inner product against the catalogue,

\[
  \mathrm{ans}_c[w] \;=\; \sum_{j=0}^{2^{\mathrm{db}}-1}
    \EvalFull(k_c)[j] \cdot D[j][w],
  \qquad c \in \{0,1\},
\]

over the $W = L / (\bits/8)$ ring words of a record ($W = 32$ at $\bits = 64$).
Because $\EvalFull(k_0)[j] - \EvalFull(k_1)[j]$ is $1$ at $j = \alpha$ and zero
elsewhere, the \emph{difference} of the two answers is exactly $D[\alpha]$.
Note the difference rather than the sum: this project omits the $(-1)^c$ factor,
per Section~\ref{sec:tech:dpf:sign}.

Two properties carry the obliviousness claim, and they are worth separating
because the second is the stronger one:

\begin{enumerate}
  \item The server touches \emph{every} record, in index order, with no
    data-dependent branch. The loop bound is the domain size and nothing inside
    the loop depends on $\alpha$.
  \item Each expanded coefficient is a \emph{share}, so it is pseudorandom at
    every index rather than zero at all but one. There is no data pattern for a
    compiler to exploit or a cache-timing observer to see, even in principle.
    This is strictly more than ``the code contains no branch''.
\end{enumerate}

The system has three servers but the DPF is $(2,2)$. We fix $\party{0}$ and
$\party{1}$ as the key holders; all three hold the catalogue, which is public
data, so replication costs nothing in privacy.

\paragraph{What it costs.} At $\mathrm{db} = 11$ a query is a single
\textbf{227-byte} key per server against a 2048-entry domain
(Table~\ref{tab:tech:keysize}), against 16~KB for the naive one-hot query the
same read would otherwise need: a $72\times$ reduction in upload. The server's
work is one $\EvalFull$ plus $2^{11} \times 32$ ring multiply-accumulates.

\paragraph{Verification.} Reads are checked byte-exactly against a cleartext
lookup: 60 random indices at $\bits = 64$ and 20 at $\bits = 128$, plus both
domain endpoints, padding slots, and rejection of a key whose domain width does
not match the server's.

One test in this group was wrong, and establishing that took more work than
accepting it would have. A natural check --- ``no word of one server's share
should equal the plaintext'' --- fails, with 18 of 32 words matching. That is
not a leak. Records pad to 256 bytes while content never exceeds 117, and the
trailing genre bytes are usually zero, so the upper words are identically zero
across the \emph{entire} catalogue; any linear combination of zeros is zero,
for every $\alpha$ equally. It is public structure, not a signal. The test now
derives from the data which words actually vary --- 14 of 32 at $\bits = 64$
and 7 of 16 at $\bits = 128$, which agree at 112 bytes and so cross-check each
other --- and asserts only on those.

\subsection{Scoring under secret sharing}
\label{sec:tech:pir:serving}

User embeddings $\userembed$ are secret; item embeddings $\itemembed$ are
public, which is \textsc{Nudge}'s design and is what makes its power iteration
cheap. The user's row is shared 2-of-3, and because $\itemembed$ is public,
$\mathbf{s} = \avec{i} \cdot \itemembed$ is a \emph{linear map applied to a
shared vector by a public matrix}. Each server computes its own share of all
$\items$ scores entirely locally: no communication, no correlated randomness,
no multiplication protocol, no rounds. This is the single largest saving in the
serving path and it comes directly from $\itemembed$ being public.

\paragraph{Scores stay at $2\fracbits$ scale.} Both factors are encoded at
$\fracbits$ fractional bits, so the product carries $2\fracbits$. Truncating
back to $\fracbits$ requires $\mathrm{Trunc}_\fracbits$, which is interactive
and belongs to the training phase, so serving decodes at $2\fracbits$ instead.
This is recorded rather than left implicit because getting it wrong scales
every score by $2^{20}$, which is \emph{monotone} and therefore completely
invisible in a ranking.

\paragraph{The seen-item mask.} The ideal functionality sets already-rated
items to $-\infty$, which a finite ring cannot represent. We use an additive
shared vector carrying $-(2^{55})$ at seen positions and zero elsewhere;
servers add it to their score shares, which is again local. Observed encoded
scores peak near $9.7 \times 10^{12}$ and the ring floor is near
$-9.2 \times 10^{18}$, so the sentinel sits about $3700\times$ below any real
score and about $256\times$ above the floor. Both margins are asserted by a
test rather than argued in prose. The comparison in the top-$\topk$ selection
is \emph{signed}: comparing these two's-complement values as unsigned would
sort masked items \emph{first}, and a smoke test inspecting only the head of
the list would not notice.

We state one thing plainly rather than overclaim it: in this phase the user
reconstructs the score vector anyway, so masking server-side is
\emph{equivalent} to masking client-side. It is implemented server-side because
that is what the ideal functionality specifies and what the composed system
will need once the servers do more than a linear map --- not because this phase
requires it.

\subsection{The float-to-ring boundary}
\label{sec:tech:pir:boundary}

Exactly one program converts a float into a ring element:
\texttt{model/export.py}. Everything downstream trusts it, so it is tested as a
contract rather than assumed. A committed file of 25 \texttt{<double> <int64>}
pairs --- zero, both signs, the quantum and half the quantum, the observed
extremes of $\userembed$, $\itemembed$ and the scores, and the two's-complement
edge --- is read by the C++ test suite, so the boundary is checked on a fresh
clone with no export run.

It found two bugs immediately, both of which would have been hard to find
later.

\textbf{Rounding mode.} \texttt{numpy.rint} and Python's \texttt{round} use
banker's rounding, half to even, so $\mathrm{rint}(0.5) = 0$; C++
\texttt{std::round} rounds half away from zero, so $\mathrm{round}(0.5) = 1$.
They disagree at every exact half-quantum. The contract test caught this on its
first run. Python was changed to match C++, since C++ is what ships.

\textbf{Quantise-then-multiply is not multiply-then-quantise.} The reference
scores were originally computed by encoding the float product; the protocol
multiplies encoded factors, which rounds $\dime$ times and then accumulates
exactly. The two disagreed on \textbf{1650 of 1682 scores}, by up to
$2.1 \times 10^{7}$ out of $4.3 \times 10^{12}$. That is about $5\times10^{-6}$
relative and irrelevant to the ranking, which is exactly why it matters: any
tolerance-based comparison would have accepted it, and the oracle would have
been quietly wrong for every later exact test. The reference now computes what
the protocol computes, in arbitrary-precision integers.

Serving is therefore checked against the oracle by \textbf{exact} equality, not
by tolerance. A linear map over a ring has no tolerance to spend, and a
tolerance here would hide precisely the endianness, sign-extension and
scale bugs the boundary is prone to. All 1682 scores match bit-for-bit.

\paragraph{Ring width, measured.} An open question in the design was whether
$\bits = 64$ suffices at MovieLens scale or whether $\bits = 128$ is needed.
Measured on the actual factorisation: $\userembed$ peaks at $61.4$,
$\itemembed$ at $0.219$ (every row has unit $L_2$ norm by construction), and at
$\fracbits = 20$ the worst-case $\dime = 16$ accumulation is
$\mathbf{47.8}$ \textbf{bits of the 63 available} --- roughly a $32{,}000\times$
margin. Quantisation costs $\max|S_q - S| = 5.3 \times 10^{-5}$, far below the
gaps between top-20 scores, so the ranking is preserved. So $\bits = 64$ is
sufficient at this scale, and the implementation remains templated on the ring
so the $\bits = 128$ arm is tested alongside it at every step.

\subsection{End to end}
\label{sec:tech:pir:demo}

\texttt{demo --user 42 --k 10} runs the whole path: the embedding is split
across three servers, each scores the catalogue locally, the user reconstructs
and selects the top ten, and each of the ten records is then fetched by a
two-server DPF-PIR read. It prints ten real film titles, and it checks itself
twice: every fetched record is compared byte-for-byte against a cleartext
lookup, and the ranking is compared position-for-position against the
independently computed Python reference. Both agree.

For user~42, masking the 112 items already rated in the training split, the top
ten are \emph{The Princess Bride}, \emph{Indiana Jones and the Last Crusade},
\emph{Toy Story}, \emph{Jerry Maguire}, \emph{Apt Pupil}, \emph{Good Will
Hunting}, \emph{Back to the Future}, \emph{The Rock}, \emph{Mission:
Impossible} and \emph{Aladdin} --- fetched at a cost of one 227-byte key per
server per title, with neither server able to say which titles they were.

\section{What an observer learns without the private read}
\label{sec:tech:leakage}

The previous section built a private read. This one measures what happens
without it, and it is the strongest result we have: it is the quantitative
answer to why the delivery layer must be private at all, and neither reference
paper reports it, because neither implements both halves of the system.

\subsection{Why a fetch is not one bit}
\label{sec:tech:leakage:why}

\nudgeplain{} publishes the item embedding matrix $\itemembed$ in the clear.
That is the design and not an oversight: each converged row of $\itemembed$ is
opened before the next is computed, which is what makes \textsc{SetOrthogonal}
a Gram--Schmidt against \emph{public} vectors and keeps the whole
power-iteration approach affordable under secret sharing.

The consequence is the one this section measures. $\itemembed$ is a complete
latent-factor model of the catalogue, known to everyone. So an adversary who
observes \emph{which} records a user fetched does not learn $j$ isolated facts;
each fetch is a constraint on the user's private embedding $\avec{i}$, a
projection onto a basis the adversary already holds. The question is how fast
those constraints collapse.

\subsection{The experiment}
\label{sec:tech:leakage:setup}

The adversary is given $\itemembed$ and the \emph{indices} of the $j$ records
the user fetched. It is given neither the scores nor $\userembed$. It computes
$\hat{a}$, the centroid of the fetched items' unit embeddings, and we report
$\cos(\hat{a}, \avec{i})$ together with the top-20 overlap between the ranking
$\hat{a}$ induces and the user's true ranking. All 943 users, at
$\dime = 16$, $\iters = 10$, against the same factorisation the demo serves.

\paragraph{The control is the load-bearing part of the design.} A
random-direction control would flatter any attack of this shape. Popular items
sit near the top of nearly everyone's ranking, so an estimator can score well
having learned only what is \emph{public}. The honest control therefore runs
the identical estimator on the $j$ globally most-popular items, ignoring the
user entirely, and the attack may claim only what it achieves \emph{above} that
line. The random control is retained to show where the floor actually is.

\begin{figure}[h]
\centering
\includegraphics[width=\linewidth]{figures/fig_attack}
\caption{Reconstruction of a user's private embedding from observed fetches
alone, over all 943 users. \textbf{Left:} cosine against the true embedding,
with the interquartile band on the attack. \textbf{Right:} the share of the
user's \emph{next} twenty recommendations the adversary can name. Both panels
share the same controls. The gap between the attack and the popularity control
is the part that is genuinely about the individual rather than about what is
popular, and it is stable at $+0.39$ to $+0.42$ across every $j$. Regenerated
by \texttt{mingw32-make figures} from
\texttt{bench/results/leakage\_attack.jsonl}.}
\label{fig:attack}
\end{figure}

\subsection{Result}
\label{sec:tech:leakage:result}

\begin{table}[h]
\centering
\small
\caption{Recovery against the number of observed fetches. The overlap column
excludes the $j$ items already seen, so it measures prediction of the user's
\emph{next} recommendations rather than repetition of the ones handed to the
adversary.}
\label{tab:attack}
\begin{tabular}{@{}rccc@{}}
\toprule
$j$ & $\cos(\hat{a}, a)$ & popularity control & next-20 overlap \\
\midrule
1  & \textbf{0.55} & 0.26 & 38\% \\
2  & 0.66 & 0.26 & 47\% \\
5  & 0.77 & 0.39 & 53\% \\
10 & \textbf{0.84} & 0.44 & \textbf{56\%} \\
20 & 0.89 & 0.47 & 55\% \\
50 & 0.92 & 0.53 & 46\% \\
\bottomrule
\end{tabular}
\end{table}

\textbf{A single observed fetch already gives $\cos = 0.55$, twice the
popularity control. Ten give $0.84$, and are enough to predict 56\% of that
user's next twenty recommendations.} The random control sits at $0.003$.

Two caveats we raise ourselves rather than wait to be asked.

First, \emph{an individual user's cosine is not evidence}. The random control's
interquartile spread is roughly $\pm 0.18$, near $1/\sqrt{\dime}$ at
$\dime = 16$, so one user scoring $0.2$ is inside the noise of guessing. Only
population means, and the margin over the popularity control, carry a claim.

Second, the overlap column is \emph{non-monotone}: it peaks near $j = 10$ and
falls by $j = 50$. That is a property of the metric, not a degradation of the
attack. The $j$ observed items are excluded from the candidate pool, so as $j$
grows the easy head of the ranking is progressively removed and what remains is
a harder tail. The cosine rises monotonically throughout, which is the check
that the estimate itself keeps improving.

\paragraph{An honest negative result.} We also implemented a second, stronger
estimator, which uses the fact that the fetched set is the \emph{top} of the
score vector and so every fetched item outranks every unfetched one --- strictly
more information than the centroid uses. It is \emph{worse} at every $j \ge 2$
(0.81 against 0.92 at $j = 50$). We report it rather than quietly dropping it,
because it sharpens the conclusion rather than weakening it: the cheap, obvious
attack is already sufficient, so an adversary here needs no sophistication.

\subsection{What it costs to close the gap}
\label{sec:tech:leakage:cost}

\begin{figure}[h]
\centering
\includegraphics[width=\linewidth]{figures/fig_pir_cost}
\caption{The private read against the two baselines a reader would otherwise
propose, at the ML-100K operating point (1682 items, 256\,B records,
$\mathrm{db} = 11$). Both panels must be read together; showing only the left
one would be advocacy rather than evaluation.}
\label{fig:pircost}
\end{figure}

Figure~\ref{fig:pircost} prices it. Against B5, downloading the entire
catalogue and filtering at home --- which is trivially private and is the first
thing a reader proposes, since 430\,KB is not obviously unreasonable ---
DPF-PIR sends \textbf{446$\times$ fewer bytes} counting both servers.

Against B1, a cleartext lookup with no privacy at all, it costs about
\textbf{50{,}000$\times$ more server time}. And read plainly,
\textbf{DPF-PIR is the slowest of the three in server CPU}, about
$5\times$ slower than B5's local copy. Its advantage is entirely in
communication.

Setting the extra computation against the saved bytes, the private read is the
better choice on any link slower than \textbf{8.0\,Gbit/s} --- which is to say
on every real network, but we would rather state the crossover than imply there
is not one. Absolute timings on this machine vary up to $4\times$ between
back-to-back runs, so these are ratios taken within a single run, as
Section~\ref{sec:tech:bench} explains.

\subsection{Would decoys do instead? No, and the reason is interesting}
\label{sec:tech:leakage:decoys}

The obvious cheap alternative to a private read is to fetch the $\topk$ real
recommendations plus $r$ decoys and let the observer sort it out. We measured
it, because ``obviously insufficient'' is not an argument.

Decoys must be drawn \emph{popularity-weighted}, not uniformly. Real
recommendations are popularity-skewed, so uniform decoys would come mostly from
the long tail and be separable by inspection alone --- an observer would
discard the unpopular ones and run the original attack on what was left.
Measuring a defence against an adversary too naive to notice that would flatter
it.

\begin{figure}[h]
\centering
\includegraphics[width=\linewidth]{figures/fig_decoy}
\caption{The decoy defence, and the floor it cannot cross. Even 40 decoys ---
five times the real traffic, at 48\,KB per query --- leave the adversary at
$\cos = 0.71$, far above the popularity control at $0.44$.}
\label{fig:decoy}
\end{figure}

\begin{table}[h]
\centering
\small
\caption{Cost and benefit of decoys, at $\topk = 10$. Bandwidth is both
servers, from Section~\ref{sec:tech:leakage:cost}.}
\label{tab:decoy}
\begin{tabular}{@{}rrrr@{}}
\toprule
decoys $r$ & fetches & $\cos(\hat{a}, a)$ & bytes/query \\
\midrule
0  & 10 & 0.843 & 9{,}660 \\
5  & 15 & 0.828 & 14{,}490 \\
10 & 20 & 0.806 & 19{,}320 \\
40 & 50 & \textbf{0.705} & 48{,}300 \\
\bottomrule
\end{tabular}
\end{table}

\textbf{The defence does not work, and the reason is the same fact that made
the popularity control necessary.} Because decoys must be popularity-weighted
to be unremarkable, they point roughly along the popularity direction --- and
the popularity direction is exactly what the control already recovers, at
$\cos = 0.44$. So decoys drag the estimate down \emph{towards} the control's
level and no further. There is a floor built into the defence, and it sits
where an adversary who watched nobody already stands.

Quintupling the traffic buys a drop of $0.14$ and leaves the adversary far
above that floor. A private read costs about a tenth of that bandwidth and
removes the observation entirely. This is a negative result about the
alternative, and it is the strongest argument for the mechanism this project
actually builds.

\medskip

Together, the parts of this section make the argument the composition exists to
make: without a private read, ten observed fetches recover the user; the cheap
alternative to a private read does not close the gap; and a private read closes
it for a few hundred bytes and a fraction of a millisecond. \nudgeplain{} leaves
this to ``other means''. This is what those means cost, and what the plausible
substitutes do not buy.

\section{On the wire}
\label{sec:tech:net}

Everything to this point is a property of the construction. This section is
about demonstrating it, because an argument that message sizes are
input-independent and a measurement that they are input-independent are
different claims, and only the second survives contact with an implementation.

\subsection{Three processes, not three objects}
\label{sec:tech:net:processes}

The servers are three separate operating-system processes. Each loads only
public data --- the item matrix $\itemembed$ and the catalogue --- and receives
one replicated share, which is useless alone. They never talk to each other:
in this slice every server-side operation is local, because $\itemembed$ is
public and $\mathbf{a}\cdot\itemembed$ is therefore a linear map on shares. The
only links are client-to-server, and that is a property of the protocol rather
than a simplification of the implementation.

Because a server serves one client to completion, nothing here is concurrent,
and \emph{no threading appears anywhere in this project}. Choosing separate
processes over threads made that true rather than merely likely.

\paragraph{Framing.} \textsc{Requirements} \S5 fixes the transport as ``plain
TCP, length-prefixed frames''. A frame is a four-byte little-endian length,
written with explicit shifts rather than a \texttt{memcpy} of a
\texttt{uint32\_t}, then exactly that many bytes.

Two details are load-bearing and both are tested. TCP is a byte \emph{stream},
so a 227-byte key handed to \texttt{send()} may arrive as 200 bytes and then
27; anything that assumes one \texttt{send()} equals one \texttt{recv()} works
on loopback with small payloads and fails on a real link, which is the worst
kind of bug because the demo passes on the development machine. The test
therefore \emph{forces} a split, by giving the listener a deliberately tiny
receive buffer, rather than hoping to observe one.

The second is that the declared length is untrusted input. It arrives over a
socket, so a peer chooses it, and reading it and then allocating that many
bytes is a four-byte denial of service. It is validated against a cap
\emph{before} it sizes anything --- the same discipline
\texttt{DpfKey::Deserialize} already applies, and the reason this parser is
covered by the hardened build alongside the other two.

\subsection{The distinguisher}
\label{sec:tech:net:distinguisher}

The phase's exit criterion originally asked for \texttt{tcpdump} on the server
links showing ``nothing but pseudorandom bytes''. That was amended, partly
because there is no \texttt{tcpdump} on the development platform, but mainly
because \emph{bytes looking random is not evidence}. Any fixed byte string
looks random if you do not know what to compare it against. The claim that
matters is that an adversary given the wire bytes cannot say which record was
fetched --- and the way to establish that is to build the adversary and measure
it losing.

\texttt{probe} records many independent queries for each of several fixed
record indices, calling $\Gen$ afresh every time, so the keys are independently
sampled; an adversary fed one key repeated many times would be measuring
nothing. Two claims are then tested separately, because they can fail
separately.

\paragraph{Length.} Every frame must be the same size whatever the index. This
follows from the construction --- key size depends only on the domain width ---
but a framing change could break it silently, so it is checked in the data.
Over \textbf{1800 recorded queries across three record indices, every frame is
228 bytes}.

\paragraph{Content.} A logistic regression is trained on the raw wire bytes and
scored on held-out data.

\begin{table}[h]
\centering
\small
\caption{Distinguishing between record indices from the wire bytes alone.
Chance is $0.5$, and it lies inside every interval. The shuffled-label control
runs the identical pipeline on deliberately destroyed labels: it agrees, which
is what shows the harness is measuring what it claims to.}
\label{tab:distinguisher}
\begin{tabular}{@{}lcc@{}}
\toprule
adversary & accuracy & 95\% CI \\
\midrule
index 0 vs 1234    & 0.490 & $[0.445,\ 0.534]$ \\
index 0 vs 777     & 0.494 & $[0.449,\ 0.538]$ \\
index 1234 vs 777  & 0.500 & $[0.455,\ 0.545]$ \\
\midrule
shuffled control   & 0.483 & $[0.439,\ 0.528]$ \\
\bottomrule
\end{tabular}
\end{table}

\begin{figure}[h]
\centering
\includegraphics[width=0.72\linewidth]{figures/fig_distinguisher}
\caption{The adversary failing. What matters is not that the points sit near
$0.5$ but that the chance line falls inside every interval; reporting a bare
``we got 49\%'' would be claiming more precision than 480 held-out trials
support.}
\label{fig:distinguisher}
\end{figure}

We state the reading carefully. This is not a proof, and a null result is not
the same as a security proof: it says a particular adversary, given a
particular feature representation, did not beat chance on this data. What it
rules out is the failure mode that actually occurs in practice --- a length or
byte-level pattern correlating with the query --- and it does so by
demonstrating an adversary losing rather than by asserting that bytes look
random.

Two things are still explicitly \emph{not} claimed. Inter-frame \emph{timing}
is not analysed; only sizes and contents are. And an observer who sees both
server links can reconstruct the record by design, which is what a $(2,2)$
scheme means and is not a weakness peculiar to this implementation.

\subsection{What it cost to say this}
\label{sec:tech:net:cost}

The networked path returns the same ten titles, in the same order, as the
in-process path, with every record still byte-exact --- which is the assertion
the demo makes rather than a claim we make about it. A full run of the
protocol, including uploading the shares and the mask and fetching ten records,
moves about 86\,KB up and 46\,KB down across the three links, framing included.

Almost all of that is the seen-item mask, which is a shared vector of length
$\items$ and therefore dwarfs the ten 227-byte keys. That is worth noticing
rather than hiding: it is an argument for masking client-side in this slice,
where it is anyway equivalent, and it is the first place a wide-area deployment
would look for a saving.

\section{What the primitive costs}
\label{sec:tech:bench}

Correctness settles whether the construction is right. This section settles
whether it is affordable, which is the question that decides if the delivery
layer is buildable at all.

\subsection{How the numbers are produced}
\label{sec:tech:bench:harness}

Every measurement is emitted as one JSON object per line, carrying the fields
fixed by the design draft: the commit hash, the host, the network profile, the
parameter tuple, the pipeline stage and phase, wall time, and bytes sent. The
commit hash, host and timestamp are filled by the emitter rather than by the
caller, so a row cannot exist without provenance. The phase is validated
against the permitted set on the way out, because a misspelled phase produces a
row that a plotting script silently drops, and that is discovered days later.

Two conventions are worth stating. Wall time is recorded \emph{per operation},
with the batch size carried alongside so batch time is recoverable. And one row
is written per repetition, never an aggregate, because the raw file is meant to
be raw and aggregation inside the producer would hide the distribution.

The benchmark writes to standard output and its human-readable table to
standard error, with the redirection living in the build system. The binary
therefore has no file handling and no assumption about the working directory.

Guarding against the compiler optimising away the work being timed is partly
structural, since the primitive lives in a separate translation unit with
link-time optimisation switched off, so its bodies are not visible to be
elided. That would stop being true the moment anyone enabled it, so every
result is additionally accumulated into a live value that is forced to memory
and printed as a checksum. Query indices are precomputed outside the timed
region, because generating them inside it would time the generator, and a
strictly increasing sweep would give the branch predictor an unrealistically
easy task on the tree path. The whole-domain output buffer is allocated once,
outside every loop, since a freshly zeroed eight megabyte buffer per call would
measure the page fault handler rather than the primitive.

\subsection{Measured cost}
\label{sec:tech:bench:results}

\begin{table}[h]
\centering
\small
\caption{Median of five repetitions, 64-bit ring unless stated, on the
development laptop with no network involved. Read the ratios rather than the
absolute figures, for the reason given below.}
\label{tab:tech:bench}
\begin{tabular}{@{}rrrrr@{}}
\toprule
domain bits & $\Gen$ ($\mu$s) & $\Eval$ ($\mu$s) &
$\EvalFull$ ($\mu$s) & $\EvalFull$, $\bits=128$ ($\mu$s) \\
\midrule
8  & 0.40 & 0.27 & 18.4   & 52.4    \\
10 & 1.40 & 0.31 & 81.9   & 199.1   \\
11 & 1.76 & 0.44 & 281.0  & 440.8   \\
12 & 1.97 & 0.47 & 379.9  & 889.5   \\
14 & 1.74 & 0.43 & 1352.7 & 4483.8  \\
16 & 1.93 & 0.53 & 5536.0 & 14343.0 \\
\bottomrule
\end{tabular}
\end{table}

Three things follow.

\emph{The delivery layer is affordable at the scale this project targets.} At
eleven domain bits, which is the padded MovieLens-100K catalogue, one server
answering one query costs a single whole-domain evaluation, measured at roughly
$0.28$~ms. Key generation on the client is under two microseconds and the key
itself is 227 bytes. Nothing here is close to being the bottleneck.

\emph{Whole-domain evaluation earns its place.} Evaluating the entire domain
through repeated single-point evaluation is between three and four times slower
than the depth-first whole-domain pass across the range measured, because the
tree walk is shared rather than repeated from the root for every index. That
ratio is the entire justification for the second entry point existing.

\emph{Key generation is not dominated by drawing randomness.} A separate
measurement of a thirty-two byte draw from the operating system generator gives
about $0.07\,\mu$s, so the two draws a key needs account for well under a
tenth of generation time even at the smallest domain. This was worth checking
rather than assuming, because if it had gone the other way the sensible
optimisation would have been buffering randomness rather than touching the
tree, and that is not where the cost is.

\subsection{A caveat about the absolute figures}
\label{sec:tech:bench:caveat}

These numbers were taken on a laptop. Running the benchmark repeatedly in quick
succession produced absolute timings differing by as much as a factor of four
between runs, while the spread across repetitions \emph{within} a single run
stayed between $1.1$ and $1.7$. That pattern is consistent with thermal
throttling rather than with measurement noise.

The consequence is stated plainly rather than smoothed over. Ratios between
operations measured in the same run are trustworthy. Absolute figures should be
read as orders of magnitude, and the conclusions above are drawn only from
ratios and from margins wide enough that a factor of four does not change them.
Every repetition is retained in the raw file with its own timestamp, so the
drift is visible to anyone who looks rather than being averaged away. Obtaining
figures that would survive being quoted precisely needs a machine that is not
thermally constrained, and that is deferred rather than faked.


\bibliographystyle{alpha}
\bibliography{references}

\end{document}
